\documentclass{osucourse}
\course{1148}

\begin{document}

\CatalogDescription
\PrerequisitesAndExclusions

\SequencingChart

\begin{Purpose}
\emph{College Algebra} provides students a college level academic experience that emphasizes the use of algebra and functions in problem solving and modeling, where solutions to problems in real-world situations are formulated, validated, and analyzed using mental, paper-and-pencil, algebraic and technology-based techniques as appropriate using a variety of mathematical notation. Students should develop a framework of problem-solving techniques (e.g., read the problem at least twice; define variables; sketch and label a diagram; list what is given; restate the question asked; identify variables and parameters; use analytical, numerical and graphical solution methods as appropriate; determine the plausibility of and interpret solutions). – Adapted from the MAA/CUPM CRAFTY 2007 College Algebra Guidelines. This course is intended to satisfy the requirements of the Ohio Board of Regents TMM001 College Algebra course with learning outcomes specified in: http://regents.ohio.gov/transfer/otm/otm-learning-outcomes.php
\end{Purpose}

\begin{Text}
\emph{College Algebra \& Trigonometry}, 1\textsuperscript{st} Edition, by Miller and Gerken, published by McGraw-Hill. ISBN: 9781259976612
\end{Text}

\begin{Technology}
All students are required to have a graphing calculator, TI-83 or TI-84. Note: Any calculators (including TI-89 and TI-92) that use a Computer Algebra System (CAS) are not permitted.
\end{Technology}

\begin{TopicsOutline}
Week 1         Section 1.7 – Inequalities

\topicsub{Section 2.3 – Functions and Relations}

Week 2         Section 2.4 – Linear Equations in Two Variables

\topicsub{Section 2.5 – Applications of Linear Equations}

Week 3         Section 9.1 – Systems of Linear Equations in Two Variables

\topicsub{Section 9.2 – Systems of Linear Equations in Three Variables}

Week 4         Section 2.6 – Transformations of Graphs

\topicsub{Section 2.7 – Analyzing Graphs of Functions}

Week 5         Test 1

\topicsub{Section 2.8 – Algebra of Functions and Composition}

Week 6         Section 3.1 – Quadratic Functions and Applications

\topicsub{Section 3.2 – Polynomial Functions}

Week 7         Section 3.3 – Division of Polynomials

\topicsub{Section 3.5 – Rational Functions}

Week 8         Section 3.5 – Rational Functions

\topicsub{Section 3.6 – Polynomial and Rational Inequalities}

Week 9         Test 2

Week 10        Section 4.1 – Inverse Functions

\topicsub{Section 4.2 – Exponential Functions}

Week 11        Section 4.2 – Exponential Functions

\topicsub{Section 4.3 – Logarithmic Functions}

Week 12        Section 4.3 – Logarithmic Functions

\topicsub{Section 4.4 – Properties of Logarithms}

Week 13        Section 4.4 – Properties of Logarithms

Week 14        Section 4.5 – Exponential and Logarithmic Equations

\topicsub{Section 4.6 – Modeling with Exponential and Logarithmic Functions}

\topicsub{Comprehensive review, Final Exam}
\end{TopicsOutline}

\end{document}
